\documentclass[11pt,letterpaper]{article}
\usepackage[margin=1in]{geometry}
\usepackage[T1]{fontenc}
\usepackage{lmodern,microtype,amsmath,amssymb,amsthm,mathtools}
\usepackage[hidelinks]{hyperref}
\setlength{\emergencystretch}{2em}
\numberwithin{equation}{section}
\newtheorem{theorem}{Theorem}[section]
\newtheorem{proposition}[theorem]{Proposition}
\newtheorem{lemma}[theorem]{Lemma}
\newtheorem{corollary}[theorem]{Corollary}
\DeclareMathOperator{\Ric}{Ric}
\DeclareMathOperator{\Rm}{Rm}
\DeclareMathOperator{\Vol}{Vol}
\DeclareMathOperator{\Crit}{Crit}
\newcommand{\R}{\mathbb R}
\title{Chiral Weyl energy and Morse heights of\\immersed homotopy four-spheres}
\author{Edward B. Baker III}
\date{Working draft, 26 November 2026}

\begin{document}
\maketitle
\begin{abstract}
A smooth Euclidean hypersurface immersion of a closed homotopy
four-sphere ties its induced Weyl energy to the index-two critical
points of its linear heights. We give sharp pointwise bounds for
each chiral Weyl block in terms of the shape operator and prove
strictness after integration on a nonempty middle-signature locus.
It follows that induced chiral energy at most \(64\pi^2/9\) recognizes
the standard smooth sphere, with no hypothesis on Yamabe sign.
Below \(8\pi^2\), a hypothetical nonstandard sphere has exactly two
middle critical points in more than three quarters of the height
directions. We prove a directional recognition criterion on a
retained immersed exterior, and give a separate test for positivity
of the induced conformal class. Examples distinguish these conditions
from first-order immersion control and from curvature control of
a local joining region. The quantitative existence and marked
height coverage needed on an arbitrary homotopy sphere remain open.
\end{abstract}

\section{Metric and immersion data}

Let \(M\) be a connected closed oriented smooth homotopy four-sphere.
We use the real chiral bundles \(\Lambda^\pm\) and the
Hilbert--Schmidt norm of their Weyl endomorphisms. For an oriented
Riemannian metric \(g\) and measurable \(P\subset M\) put
\[
 E_\pm(g;P)=\int_P |W_g^\pm|_{\mathrm{HS}}^2\,dV_g,
 \qquad E_\pm(g)=E_\pm(g;M).
\]
The chiral densities are invariant under positive conformal change,
since the Hodge splitting on two-forms is unchanged and each Weyl
endomorphism has weight \(-2\) for the change \(g\mapsto f^2g\).
On a closed homotopy sphere the signature is zero,
so \(E_+=E_-\) and the full four-index Weyl integral is \(8E_+\)
\cite[Remark 2 and (0.4)]{ChangGurskyYang2003}.

Every smooth homotopy four-sphere has an embedding in \(\R^5\).
Kervaire's Theorem~3 gives a compact smooth contractible five-manifold
with boundary \(M\) \cite{Kervaire1969HomologySpheres}. Its double is
a simply connected homology five-sphere by van Kampen and
Mayer--Vietoris, hence a homotopy five-sphere. Since \(\Theta_5=0\)
\cite{KervaireMilnor1963Groups}, the seam embeds in \(S^5\).
Removing a point away from the seam gives the claimed embedding.
The curvature and height estimates apply more generally to any
smooth hypersurface immersion on this domain.
Use the induced Riemannian metric \(g=F^*\delta\), the
orientation-compatible Gauss map \(\nu:M\to S^4\), and the
shape operator \(A=-dF^{-1}d\nu\). Write
\[
 H=\operatorname{tr}A,\quad A^\circ=A-\tfrac14 H I,\quad
 B=\left(\int_M |A^\circ|^4\,dV_g\right)^{1/2},\quad
 J=\int_M(-\det A)_+\,dV_g,\quad \omega=|S^4|=\frac{8\pi^2}{3}.
\]
Here \(s_+=\max(s,0)\). Every integral or height count uses \(M\);
distinct sheets that meet in the ambient space contribute separately.
The shape, Gauss map and heights belong to \(F\) and its induced
representative. A positive Yamabe representative of \([g]\) has the
same chiral energy. It need not itself be induced by an immersion.

\section{Sharp curvature comparisons}

In an oriented principal frame write \(k_1,\ldots,k_4\) for the
principal curvatures. The Euclidean Gauss equation gives the
sectional entries \(k_i k_j\) and
\begin{equation}\label{eq:scalar}
 R=H^2-|A|^2=\tfrac34 H^2-|A^\circ|^2.
\end{equation}
Set \(s_1=k_1k_2+k_3k_4\) and define \(s_2,s_3\) by the other
opposite-pair partitions. In the normalized self-dual and
anti-self-dual bases formed from the opposite pairs, each
diagonal curvature block has entries \(s_j/2\). Its Weyl block
subtracts \(R/12\) from each entry, and \(R=2\sum_j s_j\).
Consequently the two Weyl spectra agree, and pointwise
\begin{equation}\label{eq:weylsums}
 |W^\pm|^2=\frac1{12}\sum_{i<j}(s_i-s_j)^2 .
\end{equation}

\begin{proposition}[Shape deficit]\label{prop:deficit}
For either chirality,
\begin{equation}\label{eq:shape-deficit}
 D:=\tfrac16 |A^\circ|^4-|W^\pm|^2
 \ \ge\ \tfrac89(-\det A)_+ .
\end{equation}
If \(\det A\ge0\), equality holds exactly when the principal data
form two equal pairs. If \(\det A<0\), equality holds exactly at
positive multiples of \((-1,3,3,3)\), up to permutation and
simultaneous reversal of the signs. Both the upper bound
\(|W^\pm|^2\le |A^\circ|^4/6\) and the coefficient \(8/9\)
are pointwise sharp.
\end{proposition}
\begin{proof}
Write the trace-free principal data in the form
\[
 \tfrac12(x+y+z,\ x-y-z,\ -x+y-z,\ -x-y+z).
\]
Their squared norm is \(N=x^2+y^2+z^2\). A common shift of the
principal curvatures changes the \(s_i\) by the same amount.
For the trace-free data the first opposite-pair sum is
\(x^2-N/2\), with the corresponding expressions for \(y,z\).
Thus the Weyl eigenvalues are \((x^2-N/3)/2\) and its two
permutations, giving
\[
 D=\tfrac12(x^2y^2+x^2z^2+y^2z^2).
\]
When the determinant is nonnegative, equality is equivalent
to at most one of \(x,y,z\) being nonzero, which is the stated
two-pair condition.

For a negative determinant, reverse the normal if necessary,
order the three positive curvatures, and use homogeneity to write
the list as \((-a,1-a,1-a+v,1-a+v+u)\) with \(0<a<1\) and
\(u,v\ge0\). Let \(Z=1+2v+u\) and \(b=1-2a\).
The trace-free coordinates are \((-Z,-1-u,u-1)/2\), and
\[
 \begin{split}
 32D&=2Z^2(1+u^2)+(1-u^2)^2\ge 2Z^2+1,\\
 16(-\det A)&=(1-b^2)((Z+b)^2-u^2)\\
 &\le (2Z^2+1)(1-b^2)(1/2+b^2)
       \le \tfrac9{16}(2Z^2+1).
 \end{split}
\]
The second line uses weighted Cauchy--Schwarz for \(Z+b\),
and \((1-b^2)(1/2+b^2)=9/16-(b^2-1/4)^2\).
The first bound is an equality only at \(u=0\), since \(Z\ge1\).
Equality in the other bounds gives \(2Zb=1\) and \(b^2=1/4\).
Hence \(b=1/2\), \(Z=1\), and \(v=0\), giving the stated ratio.
The equality lists are realized at a point of a Euclidean graph
by prescribing its Hessian at zero slope.
\end{proof}

Let \(U\subset M\) be the open set of shape signature \((2,2)\).
For an immersed subdomain use the same definition on its
interior. Set \(I_2(P)=\int_{P\cap U}\det A\,dV_g\) when \(P\)
is a subdomain, and \(I_2=I_2(M)\).
\begin{proposition}[Middle-signature comparison]\label{prop:middle}
On \(U\),
\begin{equation}\label{eq:middle-weyl}
 |W^\pm|^2\ge \tfrac83\det A .
\end{equation}
The constant is pointwise sharp, with equality exactly at
\((-k,-k,k,k)\) up to permutation, \(k>0\).
Equality cannot hold on an open subset of \(U\).
Thus \(E_\pm(g;P)>\frac83 I_2(P)\) whenever \(P\) has a
nonempty open middle-signature locus.
\end{proposition}
\begin{proof}
For principal data \((-a,-b,c,d)\) with \(a,b,c,d>0\), write
\(p=abcd\). The sum \(ab+cd\) is at least \(2\sqrt p\).
The two other sums, \(-ac-bd\) and \(-ad-bc\), are at
most \(-2\sqrt p\). Two differences in \eqref{eq:weylsums}
are therefore at least \(4\sqrt p\); they yield the bound.
Equality forces \(ab=cd,\ ac=bd,\ ad=bc\), hence equal
magnitudes. A graph with this Hessian proves pointwise sharpness.

For open-set equality the smooth eigenbundles of eigenvalues
\(k,-k\) have rank two. In a rank-two eigenspace, the projection
of Codazzi for two orthogonal vector fields gives zero
directional derivative of its eigenvalue in both directions.
The two eigenspaces together imply \(dk=0\). With \(k\) constant,
for sections \(X,Y\) of opposite distributions,
\((\nabla_X A)Y=(\nabla_Y A)X\) lies in both distributions and
vanishes. Metric compatibility now makes both distributions
parallel for all derivatives. Mixed sectional curvature is zero
for such parallel distributions, in conflict with its Gauss
value \(-k^2\). The nonnegative pointwise defect cannot vanish
identically on a nonempty open set in \(U\), so its integral
there is positive.
\end{proof}
Integrating Proposition~\ref{prop:deficit} also gives
\begin{equation}\label{eq:integrateddeficit}
 E_\pm(g)+\tfrac89 J\le B^2/6 .
\end{equation}
The Weyl density can be positive where the Gauss map has
zero Jacobian. For the rank-two principal list
\((k_1,k_2,0,0)\), \eqref{eq:weylsums} gives
\(|W^\pm|^2=(k_1k_2)^2/6\).

\section{Heights and smooth recognition}

The identity \(TM\simeq\nu^*TS^4\) as oriented bundles
gives \(2\deg\nu=\chi(M)=2\). In dimension four the signed
Jacobian of \(\nu\) is \(\det A\). For almost every
\(u\in S^4\), both \(u\) and \(-u\) are regular Gauss values.
The height \(h_u=\langle F,u\rangle\) is then Morse. Its
critical points are the two inverse images, and its Hessian
there is \(\langle u,\nu\rangle g(A\cdot,\cdot)\). Reversing
the sign of this Hessian preserves index two. This genericity
asserts nondegenerate critical points. Their critical values
may coincide, including at distinct source points with the
same ambient image. The area
formula gives the Chern--Lashof total average
\cite{ChernLashof1957Curvature} and its middle-index version
\cite[\S\S2--3]{OntiVlachos2017Hypersurfaces}:
\begin{equation}\label{eq:averages}
 \frac1\omega\int_{S^4}\#\Crit(h_u)\,d\sigma
 =\frac2\omega\int_M|\det A|\,dV_g=2+\frac{4J}{\omega}.
\end{equation}
\begin{equation}\label{eq:middle-mass}
 \frac1\omega\int_{S^4}m_2(u)\,d\sigma=\frac{2 I_2}{\omega}.
\end{equation}
Here \(m_2(u)\) is the index-two count and \(d\sigma\) is
ordinary spherical area. The factor two in each formula
comes from the two normal signs.

\begin{lemma}[One middle handle]\label{lem:one}
If a Morse function on \(M\) has at most one index-two
critical point, then \(M\) is diffeomorphic to \(S^4\).
The same conclusion holds for a Morse function with at most
four critical points in total.
\end{lemma}
\begin{proof}
Reorder the handles by index, keeping their counts. The
topological argument permits small localized changes to an
auxiliary scalar function to separate repeated critical values.
It gives a marked handle presentation of the same smooth
domain. Recovering an actual regular level of the original
linear height requires its own level and marking data. The
union of the \(0\)- and \(1\)-handles is connected because
no higher handle has a disconnected attaching region.
Amalgamate the \(0\)-handles and \(1\)-handles of a spanning
tree of their incidence graph to one standard \(4\)-disk.
Every remaining \(1\)-handle reattaches by its transported
marking. From the opposite end do the same with the dual
\(0\)- and \(1\)-handles. These two operations occur in
disjoint end pieces and preserve the number of middle
handles and the actual boundary identifications. Now the
number of \(1\)- and \(3\)-handles sums to the number of
\(2\)-handles by \(\chi(M)=2\).

With no middle handle the manifold is a double of standard
disks. Each boundary diffeomorphism of \(S^3\) is isotopic
to the orthogonal map of its orientation component by
Hatcher's theorem. Extending the isotopy over a collar
and the orthogonal map over the rest of the disk identifies
the double with \(S^4\) \cite{Hatcher1983Smale}.
With one middle handle, reverse the function if needed so
that the remaining odd handle has index three. The lower
half is the trace of a framed knot in \(S^3\); its outgoing
boundary is \(S^1\times S^2\) because the upper half is a
dual genus-one \(1\)-handlebody. The first homology of
integer \(n\)-surgery is \(\mathbb Z/n\mathbb Z\), so the
framing is zero. Property R makes the knot the unknot
\cite[Corollary 8.3 and Remark 8.5]{Gabai1987FoliationsIII}.
The trace is \(S^2\times D^2\), and gluing to \(S^1\times D^3\)
by the actual boundary map gives \(S^4\) by
Laudenbach--Po\'enaru \cite[Theorem A]{LaudenbachPoenaru1972Handlebodies}.

For the second assertion, two critical points give the
double already considered. With four, the ranks \(m_j\)
satisfy \(m_j=b_j+r_j+r_{j+1}\), where \(r_j\) is the rank
of the boundary maps of the Morse chain complex. The total
rank in excess of \(b_0+b_4=2\) is two, so exactly one
nonzero \(r_j\) is one. The possibilities are one extra
adjacent-index pair. Each has at most one index-two
point and is covered by the preceding argument. For an
extra \(0,1\) pair the unique \(1\)-handle joins the distinct
disks; the dual statement handles the \(3,4\) pair.
\end{proof}

\begin{theorem}[Global immersion tests]\label{thm:recognition}
Each of the following conditions on a single immersion
recognizes the standard smooth four-sphere:
\[
 J<\omega,\qquad E_+(g)\le\frac{64\pi^2}{9},\qquad
 B\le\frac{16\sqrt2\,\pi}{3}.
\]
On a hypothetical nonstandard \(M\), an immersion with
\(E_+(g)<8\pi^2\) has \(m_2=2\) in more than three quarters
of the height directions. Four mutually orthogonal such
directions can be selected.
\end{theorem}
\begin{proof}
The total critical count is even by the Morse Euler sum.
If \(J<\omega\) its average in \eqref{eq:averages} is less
than six, so there is a Morse height with two or four critical
points. On a hypothetical nonstandard type, Lemma~\ref{lem:one}
therefore gives \(J\ge\omega\) and \(m_2(u)\ge2\) for all Morse
heights. By \eqref{eq:middle-mass}, \(I_2\ge\omega\). In particular
\(U\ne\varnothing\), and Proposition~\ref{prop:middle} gives
\(E_+>8\omega/3=64\pi^2/9\). The surplus decomposes as
\[
 E_+-\tfrac83\omega
 =E_+(g;M\setminus U)
  +\int_U\bigl(|W^+|^2-\tfrac83\det A\bigr)\,dV_g
  +\tfrac43\omega\left(\frac1\omega\int_{S^4}m_2\,d\sigma-2\right).
\]
All terms are nonnegative and the middle one is positive.
This proves inclusion of the endpoint. The pointwise
constant is sharp; the argument gives no optimal global
upper energy for recognition. If \(B\le16\sqrt2\,\pi/3\),
\eqref{eq:integrateddeficit} and \(J\ge\omega\) give
\(E_+\le64\pi^2/9\), which is impossible on the nonstandard type.

When \(E_+<8\pi^2\), \eqref{eq:middle-mass} and the
middle-signature comparison give
\(\omega^{-1}\int m_2\,d\sigma<9/4\). The excess above two
bounds the proportion of directions with \(m_2\ge3\) by
less than \(1/4\). The four first vectors of a random
orthonormal frame each have uniform marginal direction.
The union of their four failure events has probability
less than one. For each chosen direction, after the two
end reductions of Lemma~\ref{lem:one}, the pair of remaining
odd-handle counts belongs to \(\{(0,2),(1,1),(2,0)\}\).
\end{proof}

The strict endpoint for \(J\) cannot be improved by the
averaging argument alone. For \(k\ge1\), put
\(\Theta=(2k+1)\pi\), choose a nonnegative smooth \(b\) supported
in \((0,\pi/2)\) with \(C=\int b(t)\cos t\,dt>1\), and set
\[
 \rho(t)=1+b(t)+b(\Theta-t),\qquad
 r(t)=\int_0^t\rho(s)\cos s\,ds,\qquad
 z(t)=\int_0^t\rho(s)\sin s\,ds.
\]
The rotational map \((t,y)\mapsto (z(t),r(t)y)\),
\(y\in S^3\), defines an immersion on \(S^4\).
Indeed \(r(\Theta-t)=r(t)\), \(r(0)=r(\Theta)=0\),
and \(r=C+\sin t>0\) through the part between the end
bumps; monotonicity and reflection give positivity on
the remaining open arms. Each collapsing orbit is smooth
in the coordinates \(w=r y\), because the profile near
each end is a round meridian. The speed of the profile
is \(\rho\). Its normal is \((\cos t,-\sin t\,y)\) up to
a global sign. A nonaxial direction has one regular
preimage per length-\(\pi\) interval, and an axial height
has an interior critical latitude. Thus every Morse height
has \(4k+2\) critical points and \(J=k\omega\). The metric
\(\rho^2dt^2+r^2g_{S^3}\) is round-conformal: with
\(d\tau=\rho\,dt/r\) the cylinder rescaled by
\(\operatorname{sech}^2\tau\) gives the round compactification,
and the quotient of its multiplier by \(r\) extends
smoothly and positively over the round cap charts.
These immersions have zero Weyl energy and positive Yamabe
sign, including the six-critical-point case at \(J=\omega\).

\subsection{Which two-handle heights give recognition?}

\begin{proposition}[Marked surgery tests]\label{prop:two-middle-tests}
Carry the attaching tubes of one selected two-middle-point height
through the tree reductions of Lemma~\ref{lem:one}. If the odd handles lie at
one end, smooth recognition follows when the other end's two-component
link has an unknot component, has tunnel number at most two,
or contains a generalized square knot \(T_{p,q}\#T_{-p,q}\)
for coprime \(p,q\ge2\), possibly after framed slides. If there is one odd handle at each end, recognition
follows when some framed slide and order of the two middle handles
gives intermediate surgery boundary \(S^3\) or \(\#_2(S^1\times S^2)\).
It also follows if the second attaching circle on the intermediate
boundary is isotopic there to the first belt circle.
The intermediate cut belongs to the transported trace. Its identification
with a regular level of the original height would require a separate marking.
\end{proposition}
\begin{proof}
A slide inducing \(s:X'\to X\) on a trace transports the outgoing
boundary gluing \(\varphi\) to \(\varphi\circ\partial s\). At the
one-sided end surgery yields \(\#_2(S^1\times S^2)\), so the cited
unknot and tunnel-number-one conditions reduce the trace to the zero-framed unlink
\cite[Propositions 3.1--3.2]{GompfScharlemannThompson2010}. Its
opposite 1-handlebody absorbs the final marking by the extension
theorem \cite[Theorem A]{LaudenbachPoenaru1972Handlebodies}.
For tunnel number two a
split zero-framed unknot and its cancelling 3-handle give the
stabilization in Meier--Schirmer--Zupan's slide theorem
\cite[Corollary 1.3(2)]{MeierSchirmerZupan2016Trisections}. The
generalized square-knot case is the closed-manifold recognition
theorem of Meier and Zupan \cite[Theorem 1.1]{MeierZupan2022SquareKnots}.
For the mixed case apply the marked intermediate-boundary result of the
companion paper \cite{Baker2026MarkedHeights} to these transported tubes.
For an \(S^3\) cut, dual Property R identifies both halves as disks.
For the other cut, each half is
\((S^1\times D^3)\mathbin{\natural}(S^2\times D^2)\).
Surgery on the standard core sphere in one half gives \(S^1\times S^3\);
the inverse framed surgery is on a generator circle and yields \(S^4\).
The companion proves these identifications with the actual outer marking
and both circle framings. For the belt-circle isotopy, the second
handle is attached along a meridian of the first. Its framing is forced
by the outgoing \(S^1\times S^2\); the companion's belt-circle theorem
identifies the result with the double of a one-1-handle, one-2-handle
contractible manifold \(V\). The handles of \(V\times[0,1]\) cancel
in dimension five, so its boundary is standard.
\end{proof}

\section{The retained exterior and directional counts}\label{sec:exterior}

Suppose \(M=P\cup Q\) along a smooth common boundary, and
\(F:P\to\R^5\) is an immersion smooth up to the boundary. For
almost every direction \(u\), suppose \(h_u\) on \(P\) extends
to a smooth Morse function on \(M\) with at most \(k(u)\)
index-two critical points in \(\operatorname{int}Q\), where
\(k\) is nonnegative and measurable. A different scalar
extension may be used for each direction. Define
\[
 D_P=\frac12\int_{S^4}(2-k(u))_+\,d\sigma .
\]
\begin{proposition}[Exterior recognition]\label{prop:exterior}
If \(D_P>0\) and \(E_+(g;P)\le \frac83 D_P\), then
\(M\) is diffeomorphic to \(S^4\). The conclusion also holds
with \(D_P=0\) under the same energy inequality if the interior
of \(P\) has a nonempty middle-signature locus.
\end{proposition}
\begin{proof}
The Gauss image of the seam and its antipode have spherical
measure zero. At almost every remaining direction the
exterior count is at least \((2-k(u))_+\) if \(M\) is
nonstandard, by Lemma~\ref{lem:one}. The area formula restricted
to \(P\) yields \(I_2(P)\ge D_P\). If \(D_P>0\) the middle-signature
locus meets the interior of \(P\). With this locus either
inferred or supplied, Proposition~\ref{prop:middle} gives
\(E_+(g;P)>\frac83 I_2(P)\), contradicting the assumed bound.
\end{proof}

\subsection{Quadratic seam data}

Let \(Q\) consist of finitely many disjoint marked domains
\(\Omega_j\subset\R^2_x\times\R^2_y\) with smooth boundary.
Assume \(F\) extends across each marked collar by the germ
of the graph
\[
 F_j(z)=(R_j z,R_j t_j q(z)),\qquad q(z)=\tfrac12(|y|^2-|x|^2),
 \qquad R_j,t_j>0,
\]
up to a specified rigid motion of each graph. Use these
graphs on the interiors to extend the heights. Their shape
signature is \((2,2)\) by congruence with the graph Hessian.
If \(N_j\) is the image of the interior by the upward
unit normal after that rigid motion, then \(N_j\) and
\(-N_j\) are disjoint and
\begin{equation}\label{eq:normalarea}
 k(u)=\sum_j \boldsymbol{1}_{N_j\cup(-N_j)}(u),\qquad
 \sigma_j:=|N_j|=t_j^4\int_{\Omega_j}(1+t_j^2|z|^2)^{-5/2}\,dz.
\end{equation}
In fact the normal is \((-t_jD z,1)/\sqrt{1+t_j^2|z|^2}\)
before the rigid motion, with \(D=\operatorname{diag}(-1,-1,1,1)\).
Its chart \(p=u_h/u_5=-t_jD z\) is injective. The spherical
metric in this chart has one eigenvalue \((1+|p|^2)^{-2}\)
along \(p\) and three \((1+|p|^2)^{-1}\) normal to it, which
gives the area factor in \eqref{eq:normalarea}.

These seams give a middle-signature open set on the exterior
side whenever there is at least one patch: the graph has
nondegenerate shape at the seam and the exterior has the
same smooth germ there. In particular, the strict comparison
in Proposition~\ref{prop:exterior} is available. The angular
threshold is
\[
 D_P=\omega-\sum_j \sigma_j+
                \frac12\int_{S^4}(k(u)-2)_+\,d\sigma .
\]
Thus \(E_+(g;P)+\frac83\sum_j\sigma_j\le\frac83\omega\)
is a sufficient, potentially weaker criterion. For one patch
the exact threshold is \(\frac83(\omega-\sigma_1)\).
For copies of one star-shaped domain centered at the origin
with common ambient axes and horizontal marking, the \(N_j\)
are nested by the normal chart in \eqref{eq:normalarea}.
Pointwise \(\min(2,k)\) is the sum of the two largest antipodal
indicators. Hence \(D_P=\omega-\sigma_{(1)}-\sigma_{(2)}\),
where the terms are the two largest areas and an absent
term is zero. Distinct rotations require the overlap integral
in the full expression. Only the actual exterior energy and
the supplied seam data enter these recognition criteria.

\begin{proposition}[A forced middle point for a fixed graph germ]\label{prop:fixedgerm}
Let \(G:\overline\Omega\to\R^5\) be a smooth immersion agreeing
near the entire marked boundary with \(F_{R,t}\) for one
quadratic graph above. Write \(N\) for that graph's interior
one-sign normal image and \(\sigma_N=|N|\). For almost every
direction in \(N\cup(-N)\), the height of \(G\) has at least
one middle critical point in \(\Omega^\circ\). In particular
\(I_2(G;\Omega)\ge\sigma_N\) and
\(E_\pm(G^*\delta;\Omega)>\frac83\sigma_N\).
\end{proposition}
\begin{proof}
After reversing the direction if necessary, take \(u\in N\)
with both normal signs regular for \(G\). The tilted
quadratic reference has its unique critical point \(z_0\)
inside \(\Omega\), and can be written as
\(h_0=c+\frac b2(|\eta|^2-|\xi|^2)\), where \(b=Rt u_5>0\)
and \((\xi,\eta)=z-z_0\) in the negative and positive planes.
Extend \(h=\langle G,u\rangle\) smoothly by \(h_0\) outside
\(\Omega\). Choose \(L\) larger than the deviations from
\(c\) of the two heights on the patched domain, and choose
\(r_0\) with the domain inside \(|z-z_0|<r_0\) and
\(br_0^2/2>L\).

Add the same smooth radial term \(\psi(s)\),
\(s=|\xi|^2+|\eta|^2\), to both functions, with \(\psi=0\)
up to \(r_0^2\) and \(\psi'\) strictly increasing to infinity
thereafter. The new functions are proper. Every new radial
critical point has \(\eta=0\) and \(\psi'(s)=b/2\). Its value
lies below \(c-L\), since \(-bs/2+\psi(s)\) decreases until
that radius. Thus the critical points in the compact band
\([c-L,c+L]\) of the modified height are exactly those
inside the patch. The two end sublevel sets coincide with
the corresponding reference sublevel sets. For the reference
both sets retract together by contracting \(\eta\) to zero.
In the negative plane the upper set is a disk and the lower
one an annulus, from the minimum and subsequent increase of
\(-bs/2+\psi(s)\). Their relative \(H_2\) is \(\mathbb Z\).
The Morse cell attachment theorem for a compact band now
forces an index-two critical point. Equal critical values
can first be separated by small perturbations constant near
the affected critical points. The band endpoints and the
critical indices are preserved
\cite[Theorems 3.1--3.2 and Remark~3.3]{Milnor1963MorseTheory}.

Integrate over \(N\), using the two disjoint normal signs and
the area formula. This gives \(I_2(G;\Omega)\ge\sigma_N\)
and a nonempty open middle-signature locus. Proposition
\ref{prop:middle} gives the strict energy assertion.
\end{proof}
The alternating index sum from the boundary degree of the
gradient would allow an extremum in place of a middle point.
The relative sublevel group determines a specific degree.
A more general seam requires a computable fixed sublevel
pair and the auxiliary index-two count in its exterior.

\section{Yamabe sign and second-order tests}\label{sec:sign}

Write \(Q_g(f)=\int_M(6|df|^2+Rf^2)\,dV_g\),
\(\kappa=|B^4|=\pi^2/2\), and \(B_Y=32\sqrt\kappa/3=16\sqrt2\,\pi/3\).
The sharp submanifold Sobolev bound in codimension two gives
\begin{equation}\label{eq:brendle}
 4\kappa^{1/4}\|w\|_{4/3}
 \le \int_M (|dw|^2+H^2w^2)^{1/2}\,dV_g .
\end{equation}
To apply Brendle's estimate to the immersed domain, compose
\(F\) with the linear inclusion \(\R^5\subset\R^6\). The extra
normal has zero second fundamental form and the mean curvature
vector is the trace, of length \(|H|\)
\cite[Corollary 2 and \S2, Lemmas 4--6]{Brendle2021Isoperimetric}.
The multiplicity issue can be seen in the proof. For a positive
smooth weight \(w\), scale it so that the integral in
\eqref{eq:brendle} equals \(4\int w^{4/3}\).
This is possible since \(H\) is nonzero somewhere on a closed
immersion. Solve the resulting mean-zero equation
\[
 \operatorname{div}(w\nabla v)=4w^{4/3}-(|dw|^2+H^2w^2)^{1/2}.
\]
For \(|\xi|<1\), minimize \(v-\langle F,\xi\rangle\) on \(M\).
At a minimum the map \((p,y)\mapsto dF_p(\nabla v)+y\)
from the normal bundle covers \(\xi\), and
\(\nabla^2v-\langle\mathrm{II},y\rangle\) is nonnegative.
Its trace is at most \(4w^{1/3}\) by the equation and
Cauchy--Schwarz applied to the tangent and normal components
of \(\xi\). The local Jacobian is therefore at most \(w^{4/3}\).
For the shell \(\sigma<|\xi|<1\), each admissible normal
fiber in dimension two has area at most \(\pi(1-\sigma^2)\).
The area formula bounds the ambient shell volume from above
by \(\pi(1-\sigma^2)\int w^{4/3}\), with multiplicities on
the domain side. Divide by \(1-\sigma^2\) and let \(\sigma\uparrow1\).
Since \(3|B^6|=\pi\kappa\), the normalized integral
of \(w^{4/3}\) is at least \(\kappa\). Undoing the normalization
gives \eqref{eq:brendle}
on the immersed domain without an injectivity assumption.

Set \(w=|f|^3\) by smooth approximation. Cauchy--Schwarz in
\eqref{eq:brendle} gives
\[
 16\sqrt\kappa\,\|f\|_4^2\le9\|df\|_2^2+\|Hf\|_2^2.
\]
By \eqref{eq:scalar} this implies both form bounds
\begin{equation}\label{eq:form}
 Q_g(f)\ge B_Y\|f\|_4^2-\int_M(|A^\circ|^2-H^2/12)f^2\,dV_g
 \ \ge\ (B_Y-B)\|f\|_4^2+\tfrac1{12}\|Hf\|_2^2 .
\end{equation}
Consequently \(Y([g])>0\) if \(B\le B_Y\) or if
\(\|(|A^\circ|^2-H^2/12)_+\|_2<B_Y\). At \(B=B_Y\), test
on the positive first conformal-Laplacian eigenfunction.
The last term in \eqref{eq:form} is positive since a
closed Euclidean minimal immersion would have all
coordinate functions harmonic and hence constant. The
first eigenvalue has the sign of the Yamabe constant
\cite[\S\S2,11]{LeeParker1987Yamabe}.

For comparison, on a homotopy sphere a supplied nonnegative
Yamabe sign and \(E_+<8\pi^2\) already imply strict positivity.
At zero sign, the scalar-flat conformal representative
\(\tilde g\) satisfies
\[
 E_+(g)=8\pi^2+\tfrac14\|\Ric_{\tilde g}\|_2^2
\]
by Gauss--Bonnet and signature \cite[(0.4)]{ChangGurskyYang2003}.
This identity cannot determine the sign without the
nonnegativity premise.

An actual immersion with \(Y([g])>0\) and \(E_+<8\pi^2\)
enters the intrinsic minimization over positive Yamabe
classes on its fixed smooth type. For example \(B<4\sqrt3\,\pi\)
guarantees this for its induced class by \eqref{eq:form}
and \(E_+\le B^2/6\). More generally, if \(J\ge\omega\)
and positivity is available independently, the bound
\(B<4\sqrt{35}\,\pi/3\) suffices by \eqref{eq:integrateddeficit}.
With \(J<\omega\), smooth recognition already follows.
The coupled nonstandard estimate at \(B=B_Y\) gives
\(E_+\le64\pi^2/9\), and Theorem~\ref{thm:recognition}
then excludes that smooth type. The intrinsic
Chang--Gursky--Yang threshold \(E_+\le4\pi^2\) also gives
smooth recognition if \(Y>0\). The strict theorem applies in
the four-index norm, and the equality alternatives on a
nonstandard type are ruled out by simple connectivity and
\(b_2=0\) \cite[Theorems A--B]{ChangGurskyYang2003}.
The extrinsic middle-signature argument explains why the
threshold in Theorem~\ref{thm:recognition} does not require a
Yamabe condition.

\subsection{First-order control and concentration}

A regular homotopy permits changes of the second derivatives
that are invisible in a \(C^1\) limit. In a graph chart
\(x\mapsto(x,u_0(x))\), multiply its height by a cutoff to
flatten an interior cube while keeping the exterior.
Graph interpolation keeps the horizontal differential
equal to the identity throughout. In the flat cube
\([-a,a]^4\), use the perturbation
\[
 u_k(x)=\frac{\varepsilon_k}{k}\,\eta(x)
               \bigl(\cos(kx_1)+\cos(kx_2)\bigr),
 \qquad k=m\pi/a,\quad \varepsilon_k=k^{-1/2},
\]
where \(\eta=1\) near the cube and is supported inside the
flattened chart. These perturbations tend to zero in \(C^1\).
On the cube \(w_k=\sqrt{1+|Du_k|^2}\), \(dV_g=w_k\,dx\),
and the two active principal curvatures have product
\(\det D^2_{12}u_k/w_k^4\). The rank-two Weyl formula gives
\[
 E_\pm(g_k;[-a,a]^4)
 \ge\frac{4a^4\varepsilon_k^4k^4}{6(1+2\varepsilon_k^2)^{7/2}}
 \longrightarrow\infty .
\]
The determinant integral has no \(C^1\) bound either.
Using four active cosine terms of amplitude \(k^{-3/2}\),
restrict to subcubes where all cosines are bounded away
from zero and the diagonal Hessian has negative
determinant. This set has volume bounded below and the
negative determinant magnitude is of order \(k^2\).
In graph coordinates \((-\det A)_+\,dV_g=
(-\det D^2u_k)_+ (1+|Du_k|^2)^{-5/2}dx\), proving divergence.
The modifications remain in the same regular homotopy
class, relative to the outer graph. Fixed-shape patches
may also be scaled into arbitrarily small charts. The
total Weyl energy, \(J\), and the integral of
\(|A^\circ|^4\) each retain their size under ambient
homothety, since they have curvature degree four.
Shrinking support by itself supplies no control of these
integrals.

\subsection{A supported negative sign test}

There are actual immersions with negative Yamabe sign on any
of the prescribed domains. Consider the graph
\(\Sigma=(x,y,xy)\subset\R^3\), with
\[
 K=-(1+x^2+y^2)^{-2},\qquad dA=\sqrt{1+x^2+y^2}\,dx\,dy .
\]
Its parameter unit disk has curvature integral
\(-c_0=-2\pi(1-1/\sqrt2)\), and lies in the ambient-radius-two
set. Put \(\rho=|(x,y,xy)|\). Then \(|d\rho|_\Sigma\le1\) and
\(\operatorname{Area}(\rho\le R)\le16R^2\) for \(R\ge2\).
Indeed in the first quadrant integrate \(1+x+y\) over
\(0\le x\le R\), \(0\le y\le\min(R,R/x)\). For \(x\le1\)
the result is at most \(5R^2/4\). For \(x\ge1\) it is at
most \(R\log R+3R^2/2\le5R^2/2\). Sum the four quadrants.
A cutoff \(\chi\) equal to one up to \(\rho=2\), zero beyond
\(2^{N+1}\), and linear in \(\log\rho\) between the two
radii has \(\int|d\chi|^2\le64/[N(\log 2)^2]\) by summing
over the \(N\) dyadic annuli. Since \(K\le0\) and
\(\int K\chi^2\le-c_0\), choose finite \(N\) large enough
that \(6\int|d\chi|^2+2\int K\chi^2\le-c_0\).

The product graph \(\Sigma\times\R^2\subset\R^5\)
has scalar curvature \(2K\). For a nonzero compactly
supported planar test \(\varphi\) and \(f(x,y,z)=
\chi(x,y)\varphi(z/L)\) its quadratic form is at most
\[
 -c_0L^2\int_{\R^2}\varphi^2
       +6\left(\int_{\R^2}|d\varphi|^2\right)
          \left(\int_\Sigma\chi^2\,dA\right).
\]
Choose finite \(L\) to make this negative, and smooth the
cutoff inside a larger box. Cut off the graph height
outside a neighborhood of this support. A common ambient
dilation inserts the finite patch into a flattened
graph chart and rescales the supported quadratic form
by a positive constant. Extension by zero keeps its
negative sign on the closed domain. The same graph
interpolation as above preserves regular homotopy.
These examples show why quantitative selection of the
immersion and the sign test are separate from formal
immersion existence. They do not give a negative
sign conclusion for immersions in the subcritical energy band.

\section{A rotational test for global gluing budgets}\label{sec:rotational}

A rotational bridge displays the distinction between Weyl
energy and negative Gauss mass even with exact planar face
germs. Take
\[
 F(s,\theta)=(r(s)\theta,z(s)),\qquad
 r'=\cos\phi,\quad z'=\sin\phi,\qquad \theta\in S^3 ,
\]
with \(s\) arclength and \(r>0\). For
\(\nu=(-\sin\phi\,\theta,\cos\phi)\) the principal curvatures
are \(p=\phi'\) and \(q=\sin\phi/r\), the latter repeated
three times. The radial ambient acceleration is \(p\nu\),
and the orbit sphere has angular radial acceleration
\(-\theta/r\), giving \(q\). All three opposite-pair sums
are \(pq+q^2\). Thus this class has \(W^\pm=0\), and
\[
 g=ds^2+r^2 g_{S^3},\quad
 R=6q(p+q),\qquad |A^\circ|^2=\tfrac34(p-q)^2.
\]

The negative-determinant equality in
Proposition~\ref{prop:deficit} occurs for \(p=-q/3\).
At a waist with \(r=r_0>0,\ z=0,\ \phi=\pi/2\), this
ODE gives
\[
 r=r_0\sin^{-3}\phi,\qquad
 \frac{ds}{d\phi}=-3r_0\sin^{-4}\phi,\qquad 0<\phi<\pi.
\]
The first identity follows by differentiating
\(r^{1/3}\sin\phi\), and the second by the equation for \(p\).
The metric is complete since the two arclength integrals
diverge and \(r\ge r_0\). It is embedded because \(z'>0\).
This is the rotational \(H_3=0\) model
\cite[\S1.2.1, equation (7) and Theorem~2(ii)]{Palmas1999RotationHypersurfaces}.
Its determinant integral is
\[
 |S^3|\int_\pi^0\sin^3\phi\,d\phi=-\omega.
\]
The sign is negative throughout and
\(B^2=(16/3)\omega=128\pi^2/9\) by
\eqref{eq:shape-deficit} at equality.

For exact planar faces retain each arm to radius \(R\gg r_0\).
On the upper arm the outward slope parameter is
\(a(r)=(r_0/r)^{1/3}\), with \(dZ/dr=a/\sqrt{1-a^2}\).
On either arm for \(R\le r\le2R\), replace \(a\) by
\(a_\pm(r)=a(r)\chi_\pm(r/R)\), with smooth nonincreasing
cutoffs equal to one near \(R\) and zero near \(2R\).
Integrate \(a_\pm/\sqrt{1-a_\pm^2}\) from the retained
height to the new face height, and reflect the lower
arm. Each outward graph agrees with the exact parametrized
plane near its face. The planes are parallel, with
marked centers on the common axis. Only compatible
prescribed exterior germs can use these data. With \(Z\) the
ideal upper height, each absolute face height can be chosen
in \((Z(R),Z(2R))\) by moving the decrease of its cutoff
toward either end and interpolating the profiles. At fixed
\(R,r_0\) the two choices are independent. Integrating
\(dZ/dr\) shows that the separation is of order
\(R(r_0/R)^{1/3}\) for large \(R/r_0\).

Let \(\varepsilon=(r_0/R)^{1/3}\).
For the modified arm, with the global normal convention,
\(p=a_\pm'\), \(q=a_\pm/r\), up to simultaneous reversal
if the orientation is changed. The radial volume density
is \(r^3/\sqrt{1-a_\pm^2}\). Its determinant mass over the
cutoff band is exactly
\[
 |S^3|\int_0^\varepsilon \frac{v^3}{\sqrt{1-v^2}}\,dv,
\]
equal to the corresponding complete tail's mass. Thus
the marked bridge still has \(J=\omega\) and signed total
\(-\omega\). More generally, for any smooth rotational profile
with \(r>0\) going between these two outward planar ends,
the signed integral is \(-\omega\) by the primitive
\(-\cos\phi+\cos^3\phi/3\) of \(\sin^3\phi\).
It follows that its negative part costs at least \(\omega\).
For the displayed cutoffs each band's volume is \(O(R^4)\)
and \(|p|+|q|=O(\varepsilon/R)\), for a fixed cutoff family.
Consequently the flattened bridge has
\[
 W^\pm=0,\qquad B^2=128\pi^2/9+O(\varepsilon^4),\qquad
 \int R_-^2\,dV_g=O(\varepsilon^4).
\]
The shape cost is strictly larger than \(128\pi^2/9\),
because the cutoff changes the equality ratio where
\(a_\pm>0\) and \(\chi_\pm\) is decreasing. The scalar
curvature of the induced radial graph satisfies
\[
 R=\frac3{r^3}\frac{d}{dr}(r a_\pm)^2 .
\]
It is negative somewhere on each flattening arm since
the differentiated quantity goes from \(R^2\varepsilon^2\)
to zero. This sign concerns the induced representative.

The bridge is conformal to a round product, using
\(r^{-2}g=d\tau^2+g_{S^3}\), \(d\tau=ds/r\), and the
product has scalar curvature six. Its length between the
marked faces is \(L=2\log(R/r_0)+O(1)\), because the length
from the waist to \(R\) is
\(3\log\cot(\arcsin(\varepsilon)/2)\) on each side.
Suppose two supplied closed immersed sides have
positive Yamabe constants \(y_-,y_+\) and exact compatible
flat disk germs at the marked faces. Use their entire
disk exteriors and identify the parameterizations
through flat buffers. The punctured sides in their
\(r^{-2}\) conformal gauges have exact infinite round-product
ends. Sine and cosine cutoffs \(\chi_\pm\) on the joined
cylinder with \(\chi_-^2+\chi_+^2=1\) and gradients
\(O(L^{-1})\) allow \(\chi_\pm f\) to be used as compact
tests on their respective punctured sides. Conformal
covariance gives \(Q(\chi_\pm f)\ge y_\pm\|\chi_\pm f\|_4^2\)
\cite[(2.7)]{LeeParker1987Yamabe}. The localization identity is
\[
 Q(f)=\sum_\pm Q(\chi_\pm f)-
                       6\int\sum_\pm |d\chi_\pm|^2 f^2 .
\]
The cylinder has volume \(O(L)\), so the last term is at
most \(C L^{-3/2}\|f\|_4^2\) by H\"older, while
\(\sum_\pm\|\chi_\pm f\|_4^2\ge\|f\|_4^2\).
The whole joined class therefore has Yamabe constant
at least \(\min(y_-,y_+)-CL^{-3/2}\). This gives positivity
for sufficiently long bridges if the supplied signs have
a uniform positive margin. In this comparison the two
full exterior Weyl energies add, and the bridge contributes
its \(J=\omega\) and shape cost in addition to the exterior
costs. The zero Weyl energy of the bridge alone gives no
bound for the joined metric's total energy. On a joined
homotopy sphere, the signed determinant integrals of the
two full exteriors must add to \(2\omega\) to compensate the
bridge's signed integral \(-\omega\).

\section{Coverage and further challenges}

The companion on marked immersed handle constructions \cite{Baker2026ImmersedHandles}
treats framed solid-torus faces, matching bands and
conditional realizations of a local index-two attachment.
Section~\ref{sec:exterior} uses compatible germs and the energy
of a common exterior in the resulting immersion. A round
rotational bridge has spherical faces and does not supply
these solid-torus attaching maps. Local estimates do not give
the full exterior or its energy.

The companion on recognition from marked Morse heights \cite{Baker2026MarkedHeights}
treats the actual diagrams for the two-middle-handle
case. Theorem~\ref{thm:recognition} gives the angular proportion of
regular heights with exactly two middle points. The three permitted odd-handle
distributions do not themselves decide geometric cancellation.
A proof by selected heights needs a direction in the same
immersion whose actual framed levels satisfy a recognition
hypothesis. Even four orthogonal good directions from the
count give no such attaching-map information by themselves.

If a positive subcritical class passes to intrinsic fixed-type
minimization, its initial immersion need not extend to any
minimizing class. The fixed-type attainment paper treats minimization after entry
\cite{Baker2026BourAttainment}. The harmonic-annulus comparison in the marked-neck paper \cite{Baker2026MarkedNecks} takes a supplied oriented embedded comparison
with small fixed distortion, vanishing full curvature in that
gauge, harmonic or controlled-tension coordinates, and protected
placements of the two charged sides. It does not extract those
data from a Euclidean immersion, its Weyl integral or its height
counts. The rotational model above furnishes an exact local
product after conformal change, with different hypotheses on
the marked sides.

For an arbitrary homotopy sphere the remaining geometric
premise is a global immersed construction with a qualifying
energy or shape bound, or a retained immersed exterior meeting
the angular test with actual matching data. For the subcritical
height route a separate coverage mechanism must link the
available directions to the certified diagrams. For intrinsic
entry both the sign and the strict subcritical energy belong
to the same supplied conformal class. These alternatives do
not presently cover every possible smooth type.

\paragraph{Provenance.}
Prepared for Edward B. Baker III with OpenAI Codex assistance
on 25 November 2026. This standalone analysis incorporates the
parent immersion investigation and the master computations on
immersion entry, height averaging, shape deficit, rotational
joining and fixed-germ exterior recognition. The frozen
repository base was \texttt{483ad3cef5a0af4328c5d35cabbc1f5224d26b1e}.
The bounded editorial request selected GPT-6.2 Sol/ultra.
Actual serving model and reasoning effort were unavailable.
The proofs retain the hypotheses of their assessed inputs.
The setup and harmonic annulus records enter the provenance
and scope comparison, respectively.

\bibliographystyle{alpha}
\bibliography{references}
\end{document}
